% appA-tilting.tex -- Appendix A of the C3 paper (draft 2026-11-08): the second proof of Theorem 5.1(a).
% Sources: claude-20261104-c3-closure/01-PC-CLOSURE-PERIODICITY.md (c1: sections 1-2, Lemma B, repairs R1-R3);
% notes-20261104-referee-c3-closure.md (cold referee, incl. the p=2 counterexample and the fixes);
% notes-20261104-referee-c3-closure-x2.md (x2: scalar extension); Astra
% codex-20261104-runB-adjudication/03-C1-INTERFACES-AND-ERRATA.txt section 4. Harman statements checked against
% arXiv:1509.06414v3 (harman-local.txt) on 2026-11-08.

\section{The second proof of periodicity}\label{app:tilting}

This appendix proves Theorem~\ref{thm:periodicity}(a), for $m,m'>2D$.  It is the proof on
which part (a), and hence the exact degree for $n\ge2q-1$, rests; it does not use
\cite{hone2025}.  It extends natural isomorphisms only between exact functors, so the flaw in
\cite[Lemma~8.2.1]{hone2025} (Section~\ref{sec:errata}), which concerns additive functors, does
not affect it.  It transports the closure of Section~\ref{sec:periodicity-closure} along
Harman's stable equivalences~\cite{harman2015}.  Those equivalences are compatible with
induction and restriction of permutation modules; the closure needs compatibility for
\emph{arbitrary} submodules, which we prove on the Schur-algebra side with standard facts on
tilting modules (Lemma~\ref{lem:B}).

\paragraph{Status.}  T-WRITTEN.  The proof was written by a Claude agent (c1) in run
20261104-0759, reviewed by a model referee with no contact with the run (who found the
characteristic-$2$ error of Section~\ref{sec:errata} and supplied the replacement argument of
Lemma~\ref{lem:R1}) and by a Codex agent (x2, who supplied the two-way scalar extension), and
adjudicated by Astra, who accepted the coherence argument of Section~\ref{app:coherence}
conditional on the named inputs and without re-proving their foundations.  The inputs are:
Harman's Theorem~2.5 with the corrected modulus and his Corollary~1.6 and Proposition~2.7
(for $\ell=1$, which the referee checked directly); Ringel duality and the equivalence
$K^b(\mathrm{Tilt})\simeq D^b$ for highest-weight categories with finitely many
weights~\cite{ringel1991,happel1988}; the facts that tensor products of tilting modules for
$GL_N$ are tilting and that restriction to a Levi subgroup preserves tilting
modules~\cite{wang1982,donkin1998,mathieu1990}; the Hom-dimension formula for permutation
modules~\cite[13.19]{james1978}; and standard properties of the Schur
functor~\cite{green1980}.

Throughout, $k=\overline{\mathbb F}_p$ (Section~\ref{sec:periodicity-closure} reduces to this
case), $q$ is a power of $p$ with $q>D$, and $m\equiv m'\pmod q$ with $m,m'>2D$.  For a
partition $\lambda$ and $m\ge|\lambda|+\lambda_1$ put $\lambda(m)=(m-|\lambda|,\lambda)$.

\subsection{Harman's equivalences, corrected}\label{app:harman}

\paragraph{Permutation and Young modules.}  Let $\mathrm{Perm}^{\le r}_m$ be the category of
direct sums of permutation modules $M(\lambda(m))$ of $S_m$ with $|\lambda|\le r$, and
$\mathrm{Young}^{\le r}_m$ its Karoubi envelope, the direct sums of Young modules
$Y(\lambda(m))$.  Homomorphisms between permutation modules have the stable Carter--Lusztig
bases of \cite[\S2]{harman2015}; a \emph{stable map} is one whose coefficients in these bases do
not depend on $m$.

\begin{theorem}[Harman~{\cite[Thm.~2.5, Cor.~2.6]{harman2015}}, modulus corrected]\label{thm:harman25}
Let $r\ge0$ and let $q$ be a power of $p$ with $q>r$.  (For $r=0$ the categories consist of
sums of the trivial module and the equivalence is the evident one.)  For $m,m'>2r$ with $m\equiv m'\pmod q$,
sending $M(\lambda(m))$ to $M(\lambda(m'))$ and each stable map to the stable map with the
same coefficients defines an equivalence
$E^{\mathrm{P}}\colon\mathrm{Perm}^{\le r}_m\to\mathrm{Perm}^{\le r}_{m'}$, and hence an
equivalence $E^{\mathrm Y}$ of Young categories.
\end{theorem}

The structure constants of composition are binomial coefficients in $m$ of degree at most
$r$, and $\binom md\bmod p$ depends on $m$ only modulo the least power of $p$ above $d$
(Lucas).  The printed modulus $p^{\lceil\log_pr\rceil}$ is too small when $r$ is a power of
$p$ (Section~\ref{sec:errata}); with $q>r$ the proof of \cite{harman2015} applies unchanged.
For $\ell=1$, \cite[Prop.~2.7]{harman2015} states that $E^{\mathrm P}$ commutes with
restriction to $S_{m-1}$ and induction from $S_{m-1}$, up to natural isomorphism, on permutation
modules.  The relevant maps are explicit: restriction sends the basis element $h_\tau$ to
$\sum h_{\tau-e_{ij}}$ and induction adjoins a singleton block, with coefficients $0$ and $1$
that do not depend on $m$; so this can be checked directly.

\paragraph{Schur algebras.}  For $N\ge m$ let $S(N,m)$ be the Schur algebra, so that
$S(N,m)$-modules are the polynomial representations of $GL_N$ of degree $m$, and let
$\mathcal S^{\le r}(N,m)$ be the Serre subcategory generated by the simples
$L(\lambda(m)')$, $|\lambda|\le r$.  It is a highest-weight category whose tilting objects
are the $T(\lambda(m)')$~\cite[Prop.~1.3 and proof of Thm.~2.8]{harman2015}.  Let $\mathcal F$
be the Schur functor and $\tilde{\mathcal F}=\mathcal F(-)\otimes\mathrm{sgn}$; it is exact and
$\tilde{\mathcal F}(T(\lambda'))=Y(\lambda)$~\cite[Cor.~1.8(1),(2)]{harman2015}.

\begin{lemma}[full faithfulness on tilting modules]\label{lem:R1}
In every characteristic, $\tilde{\mathcal F}$ restricts to an equivalence
$\mathrm{Tilt}(\mathcal S^{\le r}(N,m))\to\mathrm{Young}^{\le r}_m$.
\end{lemma}

\begin{proof}
Harman deduces this from \cite[Cor.~1.8(4)]{harman2015}, which is false at $p=2$
(Section~\ref{sec:errata}); here is a proof that does not use it.  The exterior powers
$\Lambda^\mu V=\Lambda^{\mu_1}V\otimes\Lambda^{\mu_2}V\otimes\cdots$ are tilting, and
$T(\mu')$ is a summand of $\Lambda^\mu V$ with multiplicity one; so the tilting objects of
$\mathcal S^{\le r}(N,m)$ are the summands of the $\Lambda^{\lambda(m)}V$, $|\lambda|\le r$,
and $\tilde{\mathcal F}(\Lambda^\mu V)=M(\mu)$.

\emph{Dimensions.}  Since $\Lambda^\lambda V$ has a Weyl filtration and $\Lambda^\mu V$ a
good filtration,
$\dim\mathrm{Hom}(\Lambda^\lambda V,\Lambda^\mu V)=\sum_\nu(\Lambda^\lambda V:\Delta(\nu))
(\Lambda^\mu V:\nabla(\nu))=\sum_\rho K_{\rho\lambda}K_{\rho\mu}$, the number of
$(S_\lambda,S_\mu)$ double cosets in $S_m$, which is
$\dim\mathrm{Hom}_{S_m}(M(\lambda),M(\mu))$ in every
characteristic~\cite[13.19]{james1978}.

\emph{Faithfulness.}  $\Lambda^\mu V$ embeds in $V^{\otimes m}$ (antisymmetrisation is injective
in every characteristic).  For $N\ge m$, $V^{\otimes m}$ is projective and injective, and every
simple submodule $L$ of it has $\mathcal F(L)\ne0$~\cite[\S6]{green1980}.  If
$f\colon\Lambda^\lambda V\to\Lambda^\mu V$ is non-zero, its image is a non-zero submodule of
$V^{\otimes m}$, so it contains such an $L$, and since $\mathcal F$ is exact,
$\mathcal F(\operatorname{im}f)\ne0$, so $\tilde{\mathcal F}(f)\ne0$.

A faithful functor between Hom spaces of equal finite dimension is bijective on them.  So
$\tilde{\mathcal F}$ is fully faithful on the $\Lambda^{\lambda(m)}V$, hence on their summands;
its essential image is $\mathrm{Young}^{\le r}_m$.
\end{proof}

\begin{proposition}[Harman's Theorems~2.8 and~2.9, with the repairs]\label{prop:harman-eq}
Let $N\ge m$, $N'\ge m'$.  There is an equivalence of highest-weight categories
$\Theta\colon\mathcal S^{\le r}(N,m)\to\mathcal S^{\le r}(N',m')$ with
$\Theta(T(\lambda(m)'))\cong T(\lambda(m')')$, together with a natural isomorphism
$\zeta\colon\tilde{\mathcal F}\Theta\cong E^{\mathrm Y}\tilde{\mathcal F}$ on tilting objects.
$\Theta$ descends to an equivalence of abelian categories
$E\colon\mathcal R^{\le r}_m\to\mathcal R^{\le r}_{m'}$ between the essential images
$\mathcal R^{\le r}_m=\tilde{\mathcal F}(\mathcal S^{\le r}(N,m))\subseteq\mathrm{Rep}(S_m)$,
whose kernels, cokernels and subobjects are the ordinary ones, and $E$ restricts to
$E^{\mathrm P}$ on permutation modules up to the isomorphisms $\zeta$.
\end{proposition}

\begin{proof}
By Lemma~\ref{lem:R1} and Theorem~\ref{thm:harman25}, composing $\tilde{\mathcal F}$, then
$E^{\mathrm Y}$, then a quasi-inverse of $\tilde{\mathcal F}$ at $(N',m')$ gives an equivalence
$\mathrm{Tilt}(\mathcal S^{\le r}(N,m))\to\mathrm{Tilt}(\mathcal S^{\le r}(N',m'))$ that
preserves labels.  Ringel duality~\cite[Cor.~1.6]{harman2015} extends it to $\Theta$; on
tilting objects $\Theta$ agrees with the given equivalence up to natural isomorphism, since
both are determined by the image of a full tilting object and its endomorphism ring.  (This is
a property of the construction, not its definition; we use only this property.)

For $|\lambda|\le r<m/2$ the first part $m-|\lambda|$ of $\lambda(m)$ exceeds every other
part, so $\lambda(m)$ is $p$-singular if and only if $\lambda(m')$ is.  Hence $\Theta$ maps the
Serre subcategory generated by the $p$-singular simples, which is the kernel of
$\tilde{\mathcal F}$, onto its counterpart.  It remains to identify the Serre quotient
$\mathcal S^{\le r}/\ker\tilde{\mathcal F}$ with $\mathcal R^{\le r}_m$.  Harman applies
Gabriel's lemma to $\mathcal S^{\le r}$, which needs a right adjoint landing in
$\mathcal S^{\le r}$; that is not shown (Section~\ref{sec:errata}).  Instead: the right adjoint
$\tilde{\mathcal G}$ of $\tilde{\mathcal F}$ on all of $S(N,m)$-mod has invertible counit, so
$S(N,m)\text{-mod}/\ker\tilde{\mathcal F}\simeq\mathrm{Rep}(S_m)$ by Gabriel's
lemma~\cite{gabriel1962}; and since $\mathcal S^{\le r}$ is a Serre subcategory, for
$X\in\mathcal S^{\le r}$ and $T\le\tilde{\mathcal F}X$ the pullback
$X_T=X\times_{\tilde{\mathcal G}\tilde{\mathcal F}X}\tilde{\mathcal G}T$ is a subobject of $X$
in $\mathcal S^{\le r}$ with $\tilde{\mathcal F}X_T=T$, and lifting graphs the same way gives
every morphism as a roof inside $\mathcal S^{\le r}$.  So the quotient is the full subcategory
$\mathcal R^{\le r}_m$, closed under submodules and quotients in $\mathrm{Rep}(S_m)$.  The
exact equivalence $\Theta$ descends to an exact equivalence $E$.
\end{proof}

\subsection{Compatibility with induction and restriction}\label{app:lemmaB}

Write $E_m$ for the equivalence of Proposition~\ref{prop:harman-eq} at weight $\le D$ between
$S_m$ and $S_{m'}$, and $E'_{m-1}$ for the one at weight $\le D-1$ between $S_{m-1}$ and
$S_{m'-1}$ (note $m-1>2(D-1)$ and $m-1\equiv m'-1$).

\begin{lemma}[B]\label{lem:B}
There are natural isomorphisms
\[
  E_m\circ\mathrm{Ind}\cong\mathrm{Ind}\circ E'_{m-1}\ \text{ on }\mathcal R^{\le D-1}_{m-1},\qquad
  E'_{m-1}\circ\mathrm{Res}\cong\mathrm{Res}\circ E_m\ \text{ on }\mathcal R^{\le D-1}_{m},
\]
which on permutation modules are the isomorphisms of \cite[Prop.~2.7]{harman2015} with
$\ell=1$.
\end{lemma}

\begin{proof}
Fix $N\ge\max(m,m')+1$.  On the Schur side take
$\Phi_{\mathrm{Ind}}=-\otimes V\colon S(N,m-1)\text{-mod}\to S(N,m)\text{-mod}$ and let
$\Phi_{\mathrm{Res}}\colon S(N,m)\text{-mod}\to S(N-1,m-1)\text{-mod}$ be the part of weight
$1$ for the last factor of the Levi subgroup $GL_{N-1}\times GL_1$.

\emph{(i)}  $\tilde{\mathcal F}(M\otimes V)\cong\mathrm{Ind}\,\tilde{\mathcal F}(M)$ and
$\tilde{\mathcal F}(\Phi_{\mathrm{Res}}M)\cong\mathrm{Res}\,\tilde{\mathcal F}(M)$, naturally:
the $(1^m)$-weight space of $M\otimes V$ is $\bigoplus_i M_{(1^m)-\varepsilon_i}$, which is
the induced module, and the transposition $(m\ N)$ carries the $(1^m)$-weight space of $M$ to
the weight $(1^{m-1},0,\dots,0,1)$, $S_{m-1}$-equivariantly; the sign twist commutes with both.
\emph{(ii)}  Both functors are exact and send tilting modules to tilting
modules~\cite{wang1982,donkin1998,mathieu1990}; a weight summand of a tilting module for the
Levi subgroup is tilting.
\emph{(iii)}  By branching, $-\otimes V$ maps $\mathcal S^{\le D-1}(N,m-1)$ into
$\mathcal S^{\le D}(N,m)$ (composition factors of $\Delta(\mu)\otimes V$ lie below
$\mu+\varepsilon_i$), and $\Phi_{\mathrm{Res}}$ maps $\mathcal S^{\le D-1}(N,m)$ into
$\mathcal S^{\le D-1}(N-1,m-1)$.
\emph{(iv)}  By (i), both functors preserve the kernels of $\tilde{\mathcal F}$.

\emph{Step 1 (tilting level).}  Let $\phi\in\{\mathrm{Ind},\mathrm{Res}\}$ with Schur-side
partner $\Phi$, and $T$ a tilting object.  Then, naturally in $T$,
\[
  \tilde{\mathcal F}\Theta\Phi T\cong E^{\mathrm Y}\tilde{\mathcal F}\Phi T
  \cong E^{\mathrm Y}\phi\tilde{\mathcal F}T\cong\phi E^{\mathrm Y}\tilde{\mathcal F}T
  \cong\phi\tilde{\mathcal F}\Theta T\cong\tilde{\mathcal F}\Phi\Theta T,
\]
using $\zeta$, (i), \cite[Prop.~2.7]{harman2015} (on permutation modules, hence on their
summands), $\zeta$ again, and (i).  Both $\Theta\Phi T$ and $\Phi\Theta T$ are tilting by
(ii), and $\tilde{\mathcal F}$ is fully faithful on tilting objects (Lemma~\ref{lem:R1}), so
$\Theta\Phi\cong\Phi\Theta$ naturally on tilting objects.

\emph{Step 2 (all objects).}  The categories involved are highest-weight with finitely many
weights, so $K^b(\mathrm{Tilt})\to D^b$ is an equivalence~\cite{ringel1991,happel1988}.  An
exact functor induces a functor on $D^b$ computed termwise on tilting complexes, and the natural
isomorphism of Step~1 applied termwise gives one on $K^b(\mathrm{Tilt})$, hence on $D^b$.
Taking $H^0$ on the heart gives $\Theta\Phi\cong\Phi\Theta$ on all objects, restricting to the
isomorphism of Step~1 on tilting objects.

\emph{Step 2$'$ (restriction in weight).}  The same argument shows that $\Theta$ at weight
$\le D$, restricted to weight $\le D-1$, is isomorphic to $\Theta$ at weight $\le D-1$: both
are exact and agree with the same Young-module equivalence on tilting objects.  This is needed
because $C_t\cap F_{D-1}$ is transported at weight $D$ and then restricted.

\emph{Step 3 (descent).}  By (iv) and the universal property of Serre quotients for exact
functors and natural transformations, the isomorphism descends to $E\circ\phi\cong\phi\circ E$
on the image categories.  On permutation modules it is the isomorphism of Step~1, that is,
Harman's.
\end{proof}

\subsection{Coherence}\label{app:coherence}

Two points make the isomorphisms usable for subobjects; Astra accepted both, conditional on the
inputs above.

\emph{Independence of $N$.}  Induction uses $S(N,m-1)$ and restriction lands in
$S(N-1,m-1)$.  Truncation to the weights with last coordinate $0$ is an exact Morita
equivalence $S(N,m-1)\text{-mod}\to S(N-1,m-1)\text{-mod}$ for $N-1\ge m-1$; it preserves
tilting objects and labels and commutes with $\tilde{\mathcal F}$.  Applying Step~2 to
$\Theta^{(N)}$ and to $\Theta^{(N-1)}$ conjugated by truncation (both exact, both isomorphic
to the same Young functor on tilting objects) gives a natural isomorphism between them that is
the identity on Young modules.  So $E'_{m-1}$ does not depend on $N$, and the output of
restriction can be fed into induction.

\emph{Fixed identifications.}  Fix, on each side, the descent isomorphism
$\delta\colon E\tilde{\mathcal F}\to\tilde{\mathcal F}\Theta$ and the tilting comparison
$\alpha\colon\tilde{\mathcal F}\Theta\to E^{\mathrm Y}\tilde{\mathcal F}$ of
Proposition~\ref{prop:harman-eq}.  On tilting objects their composite descends, by full
faithfulness of $\tilde{\mathcal F}$ there (Lemma~\ref{lem:R1}), to a natural isomorphism
$\eta\colon E|_{\mathrm{Young}}\to E^{\mathrm Y}$; write $\eta'$ for the same at $m-1$.  The same
$\delta$ and $\alpha$ are used in constructing the isomorphisms of Lemma~\ref{lem:B}.  Let
$c\colon E^{\mathrm P}_m\mathrm{Ind}\to\mathrm{Ind}\,E^{\mathrm P}_{m-1}$ and
$c'\colon E^{\mathrm P}_{m-1}\mathrm{Res}\to\mathrm{Res}\,E^{\mathrm P}_m$ be Harman's
isomorphisms on permutation modules.  Unwinding the descended isomorphisms $\beta$ (induction)
and $\gamma$ (restriction), the Schur-functor isomorphisms of (i) cancel by naturality of
$\delta$ and $\alpha$, and for permutation objects $P$, $Q$
\[
  \beta_P=(\mathrm{Ind}\,\eta'_P)^{-1}\circ c_P\circ\eta_{\mathrm{Ind}\,P},\qquad
  \gamma_Q=(\mathrm{Res}\,\eta_Q)^{-1}\circ c'_Q\circ\eta'_{\mathrm{Res}\,Q}.
\]
All maps used below (the multiplication pieces $g_j$, $\varepsilon_j$, the inclusions, the
axiom maps and the constant) are stable maps defined from the same stable decompositions as
$c$.  For each inclusion $T\to P$ of a submodule into a permutation object, naturality gives a
commuting square whose ambient map is the fixed component above; taking images identifies the
transported subobject.  So in the induction below every subobject is carried to the
corresponding subobject on the smaller board, and no unknown automorphism enters.

\subsection{Transport of the closure}\label{app:transport}

\paragraph{The multiplication step as induction.}  Fix a hole $j$ and let $S_{m-1}$ fix pigeon
$m$.  Multiplication by $x_{mj}$ maps $\mathrm{Res}\,F_{D-1}(m)$ into the summand
$P_j\subseteq\mathrm{Res}\,F_D(m)$ spanned by matchings containing the edge $(m,j)$; call this
map $g_j$.  $P_j$ is a permutation module of $S_{m-1}$ of weight $\le D-1$ (matchings of size
$\le D-1$ on pigeons $[m-1]$ and holes other than $j$).  Let
$\varepsilon_j\colon\mathrm{Ind}\,P_j\to F_D(m)$ be adjoint to the inclusion.  Both are stable
maps with coefficients $0$ and $1$.  For an $S_m$-submodule $S\subseteq F_{D-1}(m)$,
\[
  \textstyle\sum_i\rho(x_{ij})S=\varepsilon_j\bigl(\mathrm{Ind}\,g_j(\mathrm{Res}\,S)\bigr),
\]
because $\sum_i\rho(x_{ij})S=\sum_{\sigma\in S_m}\sigma\,\rho(x_{mj})\,S$ and the image of
$\mathrm{Ind}\,T\to F_D$ is $S_m\cdot T$.  So every $C_t(m)$ is obtained from permutation
modules of weight $\le D$ (for $S_m$) or $\le D-1$ (for $S_{m-1}$) and stable maps by images,
preimages ($C_t\cap F_{D-1}=\iota^{-1}(C_t)$ for the inclusion $\iota$), finite sums, and
induction and restriction of submodules.  The constant is the stable map $k\to F_D(m)$.

\begin{proof}[Proof of Theorem~\ref{thm:periodicity}(a)]
By induction on $t$, $\eta\bigl(E_m(C_t(m))\bigr)=C_t(m')$ as subobjects of $F_D(m')$.  For
$t=0$, $C_0$ is the sum of the images of the axiom maps, which are stable.  For the step: $E_m$
is exact, so it preserves images, preimages and sums of subobjects, and the subobjects are the
ordinary ones (Proposition~\ref{prop:harman-eq}); stable maps go to the same stable maps up to
$\eta$ (Section~\ref{app:coherence}); induction and restriction of the subobjects
$C_t\cap F_{D-1}$ and $g_j(\mathrm{Res}\,S)$ are transported by Lemma~\ref{lem:B}, with
Step~2$'$ for the change of weight.  Both chains stabilise.  Since $E_m$ is an equivalence, the
composite $k\to F_D(m)\to F_D(m)/C_\infty(m)$ is zero if and only if its image
$k\to F_D(m')\to F_D(m')/C_\infty(m')$ is.  Lemma~\ref{lem:closure} and the reduction to
$\overline{\mathbb F}_p$ (Section~\ref{sec:periodicity-closure}) finish the proof.
\end{proof}

The strict inequality $m,m'>2D$ is used twice: in Theorem~\ref{thm:harman25} and in the
comparison of $p$-singular labels in Proposition~\ref{prop:harman-eq}.  At $m=2D$ this proof
does not apply, and Theorem~\ref{thm:periodicity}(b) rests on the conditional proof of
Section~\ref{sec:periodicity-oti}.
